\documentclass[12pt,a4paper]{article}
\usepackage[a4paper,margin=1.15in]{geometry}
\usepackage[T1]{fontenc}
\usepackage[utf8]{inputenc}
\usepackage{lmodern}
\usepackage{microtype}
\usepackage{amsmath,amssymb,amsthm,mathtools}
\usepackage{enumitem}
\usepackage[hidelinks]{hyperref}
\usepackage{titlesec}
\titleformat{\section}{\large\bfseries}{\thesection}{0.6em}{}
\titleformat{\subsection}{\normalsize\bfseries}{\thesubsection}{0.6em}{}
\newcommand{\webersection}[2]{\clearpage\section*{\S #1. #2}}
\newcommand{\booktitle}[1]{\clearpage\section*{#1}}
\newcommand{\A}{\mathfrak A}
\newcommand{\G}{\mathfrak G}
\newcommand{\Q}{\mathfrak Q}
\newcommand{\Pgr}{\mathfrak P}
\newcommand{\R}{\mathfrak R}
\newcommand{\K}{\mathfrak K}
\newcommand{\eps}{\varepsilon}
\newcommand{\zetaS}{\zeta}
\newcommand{\ii}{\mathrm i}
\newcommand{\ord}{\operatorname{ord}}
\newcommand{\disc}{\operatorname{disc}}
\newcommand{\Aut}{\operatorname{Aut}}
\newcommand{\Norm}{\operatorname{N}}
\newcommand{\Tr}{\operatorname{Tr}}
\setlength{\parindent}{0pt}
\setlength{\parskip}{0.95em}
\linespread{1.34}
\title{Heinrich Weber, \emph{Lehrbuch der Algebra}, Volume II\\English Translation: Sections 125--154\\Source pp. 466--558}
\author{}
\date{}
\begin{document}
\maketitle

\booktitle{Fifteenth Section. The Form Problem of the Group \(G_{168}\) and the Theory of Equations of the Seventh Degree}

\webersection{125}{The Resolvents of the Form Problem}

The form problem for the group \(G_{168}\), in the sense of \S 43, gives first of all, by the general theory, an equation of degree \(168\), whose coefficients are invariants.  Every divisor of the group, however, gives a resolvent of lower degree, namely of degree equal to the index of that divisor.

We shall consider here only the two most interesting cases of such divisors.  The first is the octahedral group
\[
 G_{24}=\langle \chi,\omega,\theta\rangle,
\]
and the second is the subgroup of order \(21\),
\[
 G_{21}=\langle \tau,\chi\rangle,
\]
which was already met in \S 72.  They lead respectively to resolvents of the seventh and eighth degrees.

We begin with the resolvents of degree seven.  In their construction one may start from the principal axes of the group.  A principal axis is left fixed by the group \(\chi^a\theta^b\), and under the whole group \(G_{168}\) it goes into three different lines.  Thus there must be seven triples of principal axes, and these triples are determined by an equation of the seventh degree.

Put, in the notation of \S 118,
\begin{align}
 A_1&=\alpha\gamma x_1+\gamma\beta\eps^2x_2+\alpha\beta\eps^4x_3,\notag\\
 A_2&=\beta\gamma\eps^2x_1+\alpha\beta\eps^4x_2+\alpha\gamma x_3,\tag{1}\\
 A_3&=\alpha\beta\eps^4x_1+\alpha\gamma x_2+\beta\gamma\eps^2x_3.\notag
\end{align}
Then \(A_1=0,A_2=0,A_3=0\) are the equations of three of these axes.  The substitution \(\chi\) cyclically permutes \(A_1,A_2,A_3\).  The substitution \(\theta\), when actually carried out and reduced by the formulae of \S 115, produces the corresponding interchange among the same three functions.  Since \(\omega\) is expressible by \(\theta\) and \(\chi\), it follows that the quantities \(A_1,A_2,A_3\) are permuted among themselves by the octahedral subgroup.  Their symmetric functions therefore are roots of resolvents of the seventh degree.

The simplest symmetric function of these three quantities is
\[
 A_1^2+A_2^2+A_3^2.
\]
After multiplication by a convenient numerical factor, it will be used as the unknown of the resolvent.  If the sum is arranged according to \(x_1^2+x_2^2+x_3^2\) and \(x_1x_2+x_2x_3+x_3x_1\), one obtains an expression of the form
\[
 \lambda(x_1^2+x_2^2+x_3^2)+\mu(x_2x_3+x_3x_1+x_1x_2).
\]
By substituting the values of \(\alpha,\beta,\gamma\) and simplifying with the relations among the seventh roots of unity, this assumes the form
\begin{equation}
 z=x_1^2+x_2^2+x_3^2-\frac{1-\sqrt{-7}}{2}(x_2x_3+x_3x_1+x_1x_2). \tag{2}
\end{equation}
This function is taken as one root of the resolvent of the seventh degree.

The remaining six roots are obtained from it by the substitutions \(\tau^r\).  They may all be written in a common form
\begin{equation}
 z_r=\eps^r x_1^2+\eps^{2r}x_2^2+\eps^{4r}x_3^2
 -\frac{1-\sqrt{-7}}{2}\bigl(\eps^{3r}x_2x_3+\eps^{5r}x_3x_1+\eps^{6r}x_1x_2\bigr),
 \quad r=0,1,\ldots,6. \tag{3}
\end{equation}
The coefficients of the equation whose roots are the seven quantities \(z_r\) are invariants of the group.  Their degrees are easily determined.  If \(a_\nu\) is the coefficient of the term of degree \(7-\nu\) in the unknown, after the leading coefficient has been made equal to one, then \(a_\nu\) is an invariant of degree \(2\nu\).  Since no quadratic invariant exists, one must have \(a_1=0\).  The other coefficients are expressed linearly and homogeneously by the fundamental invariants:
\[
 a_2 \text{ by } f,
\quad a_3 \text{ by } \Delta,
\quad a_4 \text{ by } f^2,
\quad a_5 \text{ by } f\Delta,
\]
\[
 a_6 \text{ by } \Delta^2, f^3,
\quad a_7 \text{ by } C,K.
\]
Thus only a finite number of numerical constants has to be computed.  They are found by comparing a few terms in the expressions of the elementary symmetric functions of the \(z_r\) with the invariant expressions already obtained in \S 123.  Apart from rational numbers these constants can contain only \(\sqrt{-7}\).

We first perform the calculation under the simplifying hypothesis \(f=0\).  This case is especially important because, like the icosahedral equation, it leads to transformation equations in the theory of elliptic functions.  It is enough to compute the first terms in the power sums \(\sum z_r^3\) and \(\sum z_r^6\), and then to apply Newton's formulae.  The last coefficient is obtained from the term \(x_1^7x_2^7x_3^7\) in the product of the seven \(z_r\).  The resulting special resolvent is
\begin{equation}
 u^7-7u^5g-7\frac{1+\sqrt{-7}}{2}u^3g^2-g^3=0. \tag{4}
\end{equation}
Here the variable has been normalized so that only the parameter \(g\) remains.

If the same substitution is made in the general resolvent, in which \(f\) is no longer set equal to zero, then one has to introduce a second parameter, say \(h\).  The coefficients of the resolvent, apart from numerical factors, occur successively in the forms
\[
 h,\\quad f,\ h,\quad hg,\quad g^2,\ h^2,\quad g^2,\ h^2g,
\]
more precisely as homogeneous linear expressions in the indicated products.  The computation is carried out exactly as in the special case.  For simplification one may put one of the variables equal to zero and still obtains sufficiently many equations for all coefficients.

The two functions \(g,h\) are fractional invariants depending only on the ratios \(x_1:x_2:x_3\).  Conversely every other invariant of the same kind is rationally expressible by \(g\) and \(h\).  For, if such an invariant is written as a quotient of two forms of equal degree without common divisor, then numerator and denominator must be integral invariants of the same degree, determined individually up to constant factors.  They cannot be of odd degree, for otherwise they would have the common factor \(K\).  Hence they are rational functions of \(f,\Delta,C\).  A monomial
\[
 f^a\Delta^bC^c
\]
of degree \(m\) satisfies
\[
 4a+6b+14c=m,
\]
and after expressing \(f\) and \(\Delta\) through the normalized parameters \(g,h\), the common power of \(C\) cancels in numerator and denominator.  Everything is therefore rational in \(g,h\).

In the same way the function \(u\) depends only on the ratios \(x_1:x_2:x_3\).  Any other function with the same property and admitting the substitutions of the subgroup \(G_{24}\) is rationally expressible by \(u,g,h\).  Indeed, by the general theorems on the form problem it is first a rational function of \(u\) and of the invariants, and in a unique way it may be written as an integral polynomial in \(u\) of degree at most six.  Since the function depends only on ratios, its coefficients do so as well, and hence are rational in \(g,h\).

\webersection{126}{Reduction of the General Seventh-Degree Resolvent to the Special One}

In the preceding section we obtained two forms of the seventh-degree resolvent of the form problem.  One of them, the special resolvent, belongs to the case \(f=0\), that is, to the case in which the required point lies on the fundamental curve.  The other belongs to an arbitrary position of the point.

We shall denote by
\[
 u_x,
\qquad g_x,
\qquad h_x
\]
the quantities \(u,g,h\) when they are formed from the point \((x)\).  The general resolvent is then written
\begin{equation}
 R(u_x,g_x,h_x)=0. \tag{1}
\end{equation}
The special resolvent is obtained from this by putting \(h_x=0\).  The quantities \(u_x,g_x,h_x\) depend only on the ratios of the coordinates of the point \((x)\).

Klein observed that the solution of the general resolvent can be reduced to the solution of the special one after adjoining the root of a biquadratic equation.  To see this, introduce beside the point \((x)\) a second point \((y)\), and form the polar of \(f\),
\begin{equation}
 f_1(x,y)=y_1f_{x_1}(x)+y_2f_{x_2}(x)+y_3f_{x_3}(x). \tag{2}
\end{equation}
If \((x)\) and \((y)\) are transformed simultaneously by the same substitution of \(G_{168}\), this polar remains unchanged.

Let the point \((x)\) be arbitrary, but require the point \((y)\) to lie both on the fundamental curve and on the polar of \((x)\).  Thus
\begin{equation}
 f(y_1,y_2,y_3)=0,
\qquad f_1(x,y)=0. \tag{3}
\end{equation}
For each point \((x)\) there are four corresponding points \((y)\).  If \((x)\) is replaced by a connected point under the group, then these four points are likewise replaced by connected points.  For the point \((y)\) the quantity \(h_y\) is zero, and \(g_y\) becomes a four-valued function of \((x)\).  The elementary symmetric functions of these four values remain invariant under \(G_{168}\).  Consequently \(g_y\) is the root of a biquadratic equation whose coefficients are rational functions of \(g_x,h_x\).  The root of this biquadratic equation has to be adjoined.

The corresponding \(u_y\) is a root of the special resolvent
\begin{equation}
 R(u_y,g_y,0)=0. \tag{4}
\end{equation}
Let the four points belonging to the same \((x)\) be denoted by
\[
 y,
\quad y',
\quad y'',
\quad y'''.
\]
For an indeterminate \(t\) form
\begin{equation}
 (t-u_y)(t-u_{y'})(t-u_{y''})(t-u_{y'''})=\Phi(t). \tag{5}
\end{equation}
This expression is left fixed by all substitutions of the subgroup \(G_{24}\), and is therefore rationally expressible by \(u_x,g_x,h_x\).

Since the line represented by the polar equation can be made arbitrary by a suitable choice of \((x)\), the point \((x)\) may be chosen in such a way that, among the four points \(y,y',y'',y'''\), no two are connected by the group.  Putting \(t=u_y\) in (5), one obtains a rational equation
\begin{equation}
 \Phi(u_y;u_x,g_x,h_x)=0. \tag{6}
\end{equation}
This equation is not satisfied after a substitution from \(G_{168}\) which changes \(u_y\) into a value different from the four values associated with the same \((x)\).  Hence, considered as an equation in \(u_x\), it has exactly one root in common with the general resolvent \(R(u_x,g_x,h_x)=0\).  If one takes the greatest common divisor of these two equations, one obtains \(u_x\) rationally through \(u_y,g_x,h_x\).  This is the desired reduction of the general resolvent to the special one.

The group \(G_{168}\) also contains the subgroup \(G_{21}\) generated by \(\chi\) and \(\tau\).  It gives rise to a resolvent of degree eight.  A root of this resolvent may be taken to be the product \(x_1x_2x_3\).  Its coefficients are invariants, of degrees rising as high as the ninety-fourth.  We shall not pursue this resolvent further here.

\webersection{127}{A Permutation Group on Seven Letters of Degree 168}

The ternary substitution group of order \(168\) has a resolvent of degree seven.  The Galois resolvent of this equation of the seventh degree is therefore of degree \(168\).  It follows that the full permutation group on seven letters, whose order is
\[
 1\cdot2\cdot3\cdot4\cdot5\cdot6\cdot7=5040=168\cdot30,
\]
contains a subgroup of order \(168\).  Kronecker first recognized the existence of this subgroup; its closer investigation is of great importance for the general theory of equations of the seventh degree.

We obtain this permutation group by examining the permutations induced by the substitutions of \(G_{168}\) among the seven quantities
\[
 z_0,z_1,\ldots,z_6
\]
of \S 125.  If two linear substitutions \(s_1,s_2\) induce the permutations \(\pi_1,\pi_2\), then their product must be read in the reverse order: the composed permutation \(\pi_1\pi_2\) is induced by \(s_2s_1\).  This is merely the convention that the substitution acts first on the variables, whereas the displayed permutation acts first on the list of symbols.  In this manner we obtain, corresponding to \(G_{168}\), a permutation group of order \(168\), denoted by \(P_{168}\).  The two groups are isomorphic.

Since \(\tau\) and \(\omega\) generate \(G_{168}\), it is enough to determine the corresponding permutations.  The substitution \(\tau\) gives the cyclic permutation
\begin{equation}
 \tau=(0,1,2,3,4,5,6). \tag{1}
\end{equation}
For \(\omega\) one uses the fact that \(z_0\) is left unchanged by the octahedral subgroup.  Thus it is enough to know the effect of \(\omega\tau^r\) on \(z_0\).  The formulae of \S 72 give the required values.  Consequently one obtains a second generating permutation, together with the cycle (1), for the group \(P_{168}\).

The same group may also be represented as the ternary linear congruence group modulo \(2\), previously studied in \S 79.  Let the seven letters be denoted by the non-zero triples of residues modulo \(2\), for instance
\begin{align*}
 z_0&=(1,0,0), & z_1&=(1,1,0), & z_2&=(1,1,1),\\
 z_3&=(0,1,1), & z_4&=(1,0,1), & z_5&=(0,1,0), & z_6&=(0,0,1).
\end{align*}
Then the substitutions corresponding to the two generators are ternary linear substitutions modulo \(2\).  In particular one may choose a matrix which leaves \(z_0,z_3,z_6\) fixed and interchanges \(z_1,z_2\) and \(z_4,z_5\).  Together with the cyclic generator it yields precisely the congruence group of order \(168\).

Both generating permutations are even.  Hence \(P_{168}\) is a subgroup of the alternating group on seven letters.

\webersection{128}{Equations of the Seventh Degree with a Group of Order 168}

We now turn to the general theory of those special equations of the seventh degree whose Galois group is reduced to \(P_{168}\).

Let
\[
 A_0,A_1,A_2,A_3,A_4,A_5,A_6
\]
be arbitrary quantities, and let
\begin{equation}
 \tau=(A_0,A_1,A_2,A_3,A_4,A_5,A_6) \tag{1}
\end{equation}
be the cyclic permutation of order seven.  If we add the transposition pair
\begin{equation}
 \omega=(A_1,A_2)(A_3,A_5), \tag{2}
\end{equation}
then the compositions and powers of these permutations generate the whole group \(P_{168}\).  The permutations corresponding to the ternary generators of \(G_{168}\) are readily derived from these.

The element \(A_0\) remains fixed by the subgroup \(P_{24}\) contained in \(P_{168}\).  It is easy to form a function of the seven quantities \(A_i\) which belongs to the group \(P_{168}\).  Take, for example, the product of the three letters left fixed by a suitable octahedral subgroup, and add the products obtained by applying the cyclic permutation.  Thus one obtains
\begin{align}
 v={}&A_0A_4A_6+A_1A_5A_0+A_2A_6A_1+A_3A_0A_2\notag\\
 &+A_4A_1A_3+A_5A_2A_4+A_6A_3A_5. \tag{3}
\end{align}
Application of the generators \(\omega\) and \(\tau\) shows that \(v\) is unchanged by all permutations of \(P_{168}\).

Since \(P_{168}\) has index \(30\) in the full symmetric group and index \(15\) in the alternating group, \(v\) is a root of an equation of degree \(30\) whose coefficients are symmetric functions of the \(A_i\), and of an equation of degree \(15\) whose coefficients also may contain the product of differences of the \(A_i\).  If the \(A_i\) are the roots of a seventh-degree equation without affect, then \(v\) is the root of a resolvent of degree \(30\).  After adjoining the square root of the discriminant, this resolvent splits into two factors of degree \(15\).

Suppose now that, besides the symmetric functions of the \(A_i\), the quantity \(v\) belongs to the rationality field.  This may hold from the beginning or may be achieved by adjoining \(v\).  Then the \(A_i\) are the roots of a special equation of the seventh degree whose Galois group is contained in \(P_{168}\).

The solution of such an equation will be obtained from the form problem of \(G_{168}\).  It remains to find functions
\[
 X_1,
\quad X_2,
\quad X_3
\]
of the roots \(A_i\) which are transformed by the permutations of \(P_{168}\) in the same way as the variables \(x_1,x_2,x_3\) are transformed by the corresponding ternary substitutions of \(G_{168}\).  If these functions have been found, then substitution of them into the invariants of \(G_{168}\) gives quantities fixed by \(P_{168}\), and hence belonging to the rationality field.  The determination of \(X_1,X_2,X_3\) from these invariant values is precisely the form problem for \(G_{168}\).

Any function of \(X_1,X_2,X_3\) changed by the substitutions of \(G_{168}\) is then a resolvent of the given equation of the seventh degree; and since both \(G_{168}\) and \(P_{168}\) are simple, it is a total resolvent.

Thus, if a non-vanishing system \(X_1,X_2,X_3\) is available, the form problem of \(G_{168}\) gives the solution of the equation of the seventh degree with group \(P_{168}\).  To reduce this solution further to the special form problem occurring in the theory of elliptic functions, namely to the case \(f=0\), one still has to solve an accessory equation of the fourth degree.  This is analogous to the square root needed to reduce the general equation of the fifth degree to the icosahedral equation.

Everything now depends on constructing the functions \(X_1,X_2,X_3\) of the roots \(A_i\).  For this purpose we shall first need a few elementary propositions from general invariant theory.

\webersection{129}{Contragredient Groups}

We have already introduced in \S 37 the notion of contragredient transformations.  Let
\[
 y=A(x),
\qquad \xi=A^{-T}(\eta)
\]
be two transposed substitutions.  The two rows of variables \((x_1,x_2,x_3)\) and \((\xi_1,\xi_2,\xi_3)\), and likewise their transformed rows, are called contragredient.

Under the substitution \(y=A(x)\) every function \(\Phi(x_1,x_2,x_3)\) goes over into a function of the \(y\)'s.  Differentiating gives the first fundamental fact:
\begin{quote}
The row of variables \((x_1,x_2,x_3)\) and the row of differential operators
\[
 \left(\frac{\partial}{\partial x_1},\frac{\partial}{\partial x_2},\frac{\partial}{\partial x_3}\right)
\]
are contragredient.
\end{quote}

By repeated application of this assertion, higher differentials transform in the corresponding symmetric powers.  In practical terms, to express the \(m\)-th derivatives with respect to \(x\) by the derivatives with respect to \(y\), one expands the relevant products in the variables and then replaces products of the contragredient variables by the corresponding differential operators.  The coefficients which occur are integral rational functions of the coefficients of the substitution.

This rule may be put in a more general form.  Suppose a form \(\varphi(x)\) is carried by the substitution \(y=A(x)\) into \(\Phi(y)\), and a second form \(\psi(\xi)\) is carried by the transposed substitution into \(\Psi(\eta)\).  Then a new transformation is obtained by replacing in \(\psi\) and \(\Psi\) the powers and products of the variables by the corresponding differentiations of \(\varphi\) and \(\Phi\).  In other words, one may replace symbolic products in the contragredient variables by differential operators acting on the original form.

From the composition law of linear substitutions follows immediately:
\begin{quote}
If \(A\) runs through a group \(G\), then the transposed inverse substitution \(A^{-T}\) runs through a group \(G_1\), called contragredient to \(G\).  If \(A,B\in G\) correspond to \(A_1,B_1\in G_1\), then the order of multiplication is reversed.  The two groups are isomorphic if to \(A\) one assigns not \(A_1\), but \(A_1^{-1}\).
\end{quote}
The invariants of \(G_1\) are called the contravariants of \(G\).  Conversely, the invariants of \(G\) are the contravariants of \(G_1\).

The differential rule just stated gives two useful principles.

First, if in a contravariant of \(G\) the powers and products of the variables are replaced by the corresponding derivatives of an invariant of \(G\), the result is again an invariant of \(G\).

Second, if in an invariant of \(G\) the powers and products of the variables are replaced by the corresponding derivatives of a contravariant, then the result is again a contravariant.

These elementary rules are the mechanism by which the required functions for the seventh-degree equation will be formed.

\webersection{130}{Solution of the Seventh-Degree Equation with Group \(P_{168}\) by the Form Problem of \(G_{168}\)}

The ternary linear substitution group \(G_{168}\) has the remarkable property of being contragredient to itself.  Its generating substitutions \(\tau\) and \(\omega\) are unchanged by transposition; hence the transpose of any substitution of the group is again a substitution occurring in the group.  Thus the contravariants of \(G_{168}\), apart from the names of the variables, coincide with its invariants.

The octahedral subgroup contained in \(G_{168}\), however, is not contragredient to itself as a subgroup.  For example, the transpose of one of its substitutions is a substitution of \(G_{168}\) which need not lie in the octahedral subgroup.  This distinction is important for what follows.

Let \(\eta_0\) be any function of the seven quantities
\[
 A_0,A_1,\ldots,A_6
\]
belonging to the subgroup \(P_{24}\); the root \(A_0\) itself may be taken as an example.  By the cyclic permutations \(\tau^r\) it goes into
\[
 \eta_0,\eta_1,\ldots,\eta_6.
\]
The sums
\begin{equation}
 \omega_r=\sum_{s=0}^{6}\eps^{rs}\eta_s,
 \qquad r=0,1,\ldots,6, \tag{1}
\end{equation}
give a system of functions which are linearly, but not ternarily, transformed by the permutations of \(P_{168}\).  This is because the \(\eta_s\) may be expressed linearly in terms of the \(\omega_r\).

A ternarily transforming system of three functions is obtained as follows.  Take three functions \(\eta_0,\eta'_0,\eta''_0\), each belonging to the subgroup \(P_{24}\), and form from them the corresponding Fourier sums
\[
 \omega_r,
\quad \omega'_r,
\quad \omega''_r.
\]
Using the contragredient invariants and the differential rule of \S 129, one constructs a trilinear expression in these three rows which remains invariant when a permutation from \(P_{168}\) is applied to the roots and the corresponding contragredient substitution is applied to the variables.

Let \(f(\xi_1,\xi_2,\xi_3)\) denote the biquadratic contravariant
\begin{equation}
 f(\xi_1,\xi_2,\xi_3)=\xi_1^3\xi_2+\xi_2^3\xi_3+\xi_3^3\xi_1, \tag{2}
\end{equation}
written in variables contragredient to \((x_1,x_2,x_3)\).  If in this contravariant the products of the variables are replaced by the corresponding derivatives of a suitable invariant expression in the three rows \(\omega,\omega',\omega''\), the result is a linear form
\begin{equation}
 L=X_1\xi_1+X_2\xi_2+X_3\xi_3. \tag{3}
\end{equation}
The coefficients \(X_1,X_2,X_3\) are functions of the roots \(A_i\).  The construction is such that \(L\) remains invariant if the roots \(A_i\) are subjected to a permutation \(\pi\) of \(P_{168}\) and simultaneously the variables \(\xi\) are subjected to the corresponding contragredient substitution.

If \(\pi\) carries \(X_1,X_2,X_3\) into \(Y_1,Y_2,Y_3\), and if
\[
 A=\begin{pmatrix}
 a_1&a_2&a_3\\ b_1&b_2&b_3\\ c_1&c_2&c_3
 \end{pmatrix}
\]
is the corresponding ternary substitution, invariance of \(L\) gives
\begin{equation}
 Y_1\eta_1+Y_2\eta_2+Y_3\eta_3=X_1\xi_1+X_2\xi_2+X_3\xi_3. \tag{4}
\end{equation}
Expanding this identity yields
\begin{align}
 Y_1&=a_1X_1+b_1X_2+c_1X_3,\notag\\
 Y_2&=a_2X_1+b_2X_2+c_2X_3,\tag{5}\\
 Y_3&=a_3X_1+b_3X_2+c_3X_3.\notag
\end{align}
Thus application of the permutation \(\pi\) to the functions \(X_1,X_2,X_3\) has exactly the same effect as the corresponding linear substitution of \(G_{168}\).  These are therefore precisely the functions required at the end of \S 128.

\webersection{131}{Possibility of Determining the Functions \(X_1,X_2,X_3\)}

To make the reduction of the equation of the seventh degree with group \(P_{168}\) to the form problem of \(G_{168}\) complete, one point remains to be checked.  We have to show that the initial functions \(\eta_0,\eta'_0,\eta''_0\) can be chosen in such a way that the functions \(X_1,X_2,X_3\) do not vanish identically.

For this purpose the construction of the quantities \(X_i\) has to be examined a little more closely.  When the substitutions described in the preceding section are inserted into the symbolic formulae, the coefficients can be written as trilinear forms in the three rows
\[
 \eta_r,
\quad \eta'_r,
\quad \eta''_r.
\]
Their coefficients are rational functions of the seventh root of unity \(\eps\).  In a convenient abbreviated determinant notation the three functions take the form
\begin{align}
 X_1&=(q_3,p_2,q_1)+(p_1,q_1,q_2)+(p_1,q_3,p_3),\notag\\
 X_2&=(q_2,q_1,p_3)+(q_3,p_2,q_2)+(q_2,p_2,q_1),\tag{1}\\
 X_3&=(q_2,q_3,q_2)+(q_3,q_1,p_3)+(q_2,p_2,p_3).\notag
\end{align}
The precise letters are less important than the consequence: these are not identically zero trilinear forms.

Indeed the variables may be linearly replaced by Fourier coordinates.  Put
\begin{equation}
 \eta_r=\sum_{s=0}^{6}a_s\eps^{rs},
 \qquad
 \eta'_r=\sum_{s=0}^{6}a'_s\eps^{rs},
 \qquad
 \eta''_r=\sum_{s=0}^{6}a''_s\eps^{rs}. \tag{2}
\end{equation}
The determinant of this transformation is, up to sign, the product of all differences \(\eps^r-\eps^s\), and is not zero.  Hence the transformation is invertible.  If \(X_1,X_2,X_3\) were identically zero in the \(\eta\)'s, they would also be identically zero in the \(a\)'s.  But the explicit formulae show that suitable values of the \(a\)'s make the three functions non-zero.

Since the coefficients of these forms lie in the cyclotomic field of seventh roots of unity, one may choose rational numerical values of the auxiliary variables so that the functions \(X_1,X_2,X_3\) assume prescribed non-zero values.  The desired system of functions therefore exists.  This completes the proof that the solution of the special seventh-degree equations with group \(P_{168}\) is obtained by the form problem of \(G_{168}\), with only the auxiliary equation of degree four required in passing from the general to the special form of that problem.

\booktitle{Fourth Book. Algebraic Numbers}
\booktitle{Sixteenth Section. Numbers and Functionals of an Algebraic Field}

\webersection{132}{Definition of Algebraic Numbers}

A number \(\theta\) is called an algebraic number if it satisfies an algebraic equation with rational coefficients,
\begin{equation}
 f(\theta)=0. \tag{1}
\end{equation}
Among all such equations there is one of least degree.  Let it be
\begin{equation}
 f(x)=x^n+a_1x^{n-1}+\cdots+a_n=0, \tag{2}
\end{equation}
chosen monic.  This equation is irreducible over the field of rational numbers.  For if \(f(x)\) decomposed into two rational factors of smaller degree, then \(\theta\) would satisfy one of those factors, contrary to the minimal choice of \(n\).

If \(n\) is the degree of the rational equation of least degree satisfied by \(\theta\), then \(\theta\) is called an algebraic number of degree \(n\).

\webersection{133}{Integral Algebraic Numbers}

An algebraic number \(\theta\) is called an integral algebraic number, or simply an integer in the algebraic sense, if it satisfies an equation
\begin{equation}
 \theta^m+A_1\theta^{m-1}+\cdots+A_m=0 \tag{1}
\end{equation}
whose coefficients \(A_i\) are rational integers.

It is enough that among the infinitely many equations of this form satisfied by \(\theta\) there be one with integral coefficients.  The integral algebraic numbers contain, as a special case, the ordinary rational integers.  Positive rational integers will also be called natural numbers.  When we speak of integers without qualification in this part, we mean algebraic integers, rational or irrational.

First, if an algebraic integer is rational, it is necessarily a rational integer.  For let \(\theta=P/Q\), with \(P,Q\) rational integers without common divisor and \(Q>0\).  Substitution in (1) gives
\[
 P^m+A_1P^{m-1}Q+\cdots+A_mQ^m=0.
\]
Every prime divisor of \(Q\) would then have to divide \(P\), which is impossible unless \(Q=1\).

Second, the sum, difference, and product of two algebraic integers are again algebraic integers.  Let \(\alpha,\beta\) be algebraic integers of degrees \(\mu,\nu\).  If
\[
 \omega=\alpha+\beta,
\quad \alpha-\beta,
\quad \alpha\beta,
\]
then the \(\mu\nu\) quantities obtained from the conjugates of \(\alpha\) and \(\beta\) may be used as a finite system closed under multiplication by \(\omega\), with integral coefficients after reduction by the equations of \(\alpha\) and \(\beta\).  Eliminating this system gives a monic equation with integral coefficients for \(\omega\).  Hence \(\omega\) is an algebraic integer.

Third, if \(f(x)\) is irreducible over the rational numbers and one root of \(f(x)=0\) is an algebraic integer, then all roots are algebraic integers.  Indeed any rational polynomial vanishing at one root is divisible by the irreducible polynomial \(f\), and therefore vanishes at every conjugate root.  Applying an integral equation for the first root gives the same assertion for all conjugates.  The elementary symmetric functions of the roots then show that the coefficients of the monic irreducible equation are rational integers.

Thus, if \(\theta\) is an algebraic integer, then its equation of least degree, written monically, has integral coefficients.

Finally, every algebraic number can be made integral by multiplication by a natural number.  If
\[
 \theta^m+A_1\theta^{m-1}+\cdots+A_m=0
\]
and the rational coefficients have common denominator \(a\), then multiplication by \(a^m\) gives a monic equation with integral coefficients for \(a\theta\).  Thus \(a\theta\) is an algebraic integer.

\webersection{134}{Algebraic Fields}

In the thirteenth section of the first volume we saw how, from a root \(\theta\) of an equation of degree \(n\) irreducible over a field, one obtains an algebraic field \(\R(\theta)\) over that field.  The fields derived from the various roots of the irreducible equation are called conjugate fields; some of them, or even all of them, may coincide.

Let \(R\) denote the rational field.  Every algebraic number \(\theta\) of degree \(n\) gives rise to an algebraic number field
\[
 R(\theta),
\]
which will be called briefly an algebraic number field of degree \(n\).  It has been shown earlier that one can always find an algebraic number field containing any finite number of prescribed algebraic numbers.  Hence no generality is lost, in any discussion involving finitely many algebraic numbers, by supposing that they all lie in one algebraic number field.

Every number \(\omega\) of such a field can be represented as an integral function \(\varphi(\theta)\) of \(\theta\) with rational coefficients.  To \(\omega\) there correspond in the conjugate fields certain conjugate numbers.  These may coincide in part; accordingly one distinguishes primitive and imprimitive numbers of the field.  A symmetric function of the conjugates is rational.

Among the symmetric functions of the conjugates two are especially important: their sum and their product.  The first is called the trace, the second the norm of \(\omega\).  We write
\[
 S(\omega),
\qquad N(\omega).
\]
Repeated conjugates are counted with their multiplicities.

Since the four fundamental operations can be carried out in every algebraic number field as in the rational field, one asks how far the elementary laws of arithmetic familiar from rational integers remain valid in an arbitrary algebraic number field.  The first question concerns factorization of integers into prime factors.  In general this factorization fails if one uses only the numbers of the algebraic field.  The computational material must be enlarged in order to restore simple laws.

Kummer first solved this great problem for numbers built from roots of unity, the cyclotomic numbers, by creating ideal numbers.  Dedekind gave another general solution, without exceptions, by introducing the ideals named after him.  Kronecker gave a third approach through an arithmetical theory of algebraic quantities.  Weber will use here a method close to these theories but adapted to the applications that follow.

\webersection{135}{Integral Functions in an Algebraic Field}

Let \(\Omega\) be an algebraic number field.  We shall consider integral functions in any number of independent variables whose coefficients are numbers of \(\Omega\).  Such a function is a finite sum of terms
\[
 \alpha x^r y^s z^t\cdots,
\]
where the exponents are non-negative integers and \(\alpha\in\Omega\).  We write
\begin{equation}
 \varphi(x,y,z,\ldots)=\sum \alpha x^r y^s z^t\cdots . \tag{1}
\end{equation}
The expression is always supposed to be collected, so that the same combination of exponents does not occur twice.  Two such functions are equal only when they contain the same monomials with the same coefficients.  An integral function is zero only when all its coefficients are zero.

Sums, differences, and products of integral functions are again integral functions.  Moreover, rational numerical values may be assigned to the variables so that any prescribed finite set of non-zero such functions remains non-zero.

Integral functions over the rational field occur as a special case:
\[
 \Phi(x,y,z,\ldots)=\sum a x^r y^s z^t\cdots,
\]
with rational coefficients.  If the coefficients are rational integers without a common divisor, the function is called primitive.  The product of two primitive integral functions is again primitive.  More generally, the divisor or content of a product is the product of the contents of its factors.

If the coefficients of an integral function in \(\Omega\) depend on an algebraic number \(\theta\), the conjugate functions are obtained by replacing \(\theta\) by its conjugates.  A function belongs to the rational field exactly when all of its conjugates are equal.  Its norm is the product of all conjugate functions,
\[
 N(\varphi)=\varphi\varphi'\varphi''\cdots,
\]
and is an integral function over the rational field.  It is divisible by \(\varphi\), and its quotient is the product of the remaining conjugates.

These elementary facts will be used repeatedly when functions are treated as arithmetic objects.

\webersection{136}{The Functionals of an Algebraic Field \(\Omega\) and the Extended Field \(\Omega^*\)}

The variables occurring in the theory of algebraic numbers do not have the meaning of variables ranging over independently varying numerical systems, as in function theory.  They are only symbols of calculation.  What matters is the system of coefficients; the variables serve to make the familiar rules of literal calculation available for that system.  It is not excluded that one sometimes substitutes numerical values for them, but this is not their primary role.

We therefore introduce integral and fractional functions in any number of variables, with coefficients in the algebraic field \(\Omega\), and agree to calculate with them according to the ordinary rules of algebra.

Every such expression \(\omega\) can be written as a quotient of two integral functions in \(\Omega\),
\[
 \omega=\frac{\varphi}{\psi},
\qquad \psi\ne0.
\]
Two quotients are equal if \(\varphi\psi_1=\varphi_1\psi\).  If numerator and denominator have no common divisor, the quotient is irreducible.  Among all representations there is one in irreducible form; in it numerator and denominator are determined by \(\omega\) up to a common factor belonging to \(\Omega\).

Such a fractional expression will be called a \emph{functional} of the field \(\Omega\).  Integral functions and numbers themselves are special cases.  The four fundamental operations are applicable to the functionals in the same extent as to numbers.  Thus the totality of all functionals of \(\Omega\) is again a field, denoted by \(\Omega^*\), the functional field of \(\Omega\).  The number field \(\Omega\) is contained in \(\Omega^*\).

A functional whose coefficients are rational numbers is a rational functional.  The rational functionals form the functional field \(R^*\) of the rational field \(R\), and \(R^*\) is contained in every algebraic functional field \(\Omega^*\).

Every rational functional can be represented as a quotient of two integral rational functions with integral coefficients.  If a rational integral function has integral coefficients with greatest common divisor \(a\), it can be written
\[
 aE,
\]
where \(E\) is primitive.  For a rational functional \(A\) write
\begin{equation}
 A=aE. \tag{1}
\end{equation}
The natural number \(a\) is called the absolute value of \(A\); the primitive factor \(E\) is called a rational unit.  In the special case in which \(A\) is a number, \(a\) is the ordinary absolute value and \(E=\pm1\) is its sign.  The general definition may look unfamiliar at first, but it is the right analogue for the calculations to follow.

The absolute value of a product of rational functionals is the product of their absolute values.  If \(\omega\) is a functional in \(\Omega^*\), then its norm \(N(\omega)\) is rational.  The absolute value of this rational functional is called the absolute norm of \(\omega\), and is denoted by
\begin{equation}
 N(\omega)=N_a(\omega)E(\omega). \tag{2}
\end{equation}
Here \(N_a(\omega)\) is a positive rational number and \(E(\omega)\) a rational unit.  The absolute norm is multiplicative:
\begin{equation}
 N_a(\alpha\beta)=N_a(\alpha)N_a(\beta). \tag{3}
\end{equation}

Replacing all coefficients of a functional by the corresponding numbers in a conjugate field gives a conjugate functional.  Every identity between functionals remains valid after this simultaneous replacement.  If \(t\) is a new variable not occurring in \(\omega\), the function
\begin{equation}
 \Phi(t)=N(t-
\omega)=t^m+A_1t^{m-1}+\cdots+A_m \tag{4}
\end{equation}
has rational functional coefficients and vanishes for \(t=\omega\).  More generally, there may be equations of lower degree with rational functional coefficients that vanish at \(\omega\).

A functional \(\omega\) is called integral if it satisfies a monic equation
\begin{equation}
 \omega^m+A_1\omega^{m-1}+\cdots+A_m=0 \tag{5}
\end{equation}
with coefficients \(A_i\) integral rational functionals.  This definition extends the definition of algebraic integers from numbers to functionals.

\webersection{137}{Integral Functionals}

The integral functionals enjoy the same elementary closure properties as the algebraic integers.

First, if two functionals are integral, then their sum, difference, and product are integral.  This follows by the same determinant argument used for algebraic integers: the products of a suitable finite system by the new functional are expressible linearly in that system, with integral rational functional coefficients, and elimination gives a monic integral equation.

Second, if \(\omega\) is integral and \(\Phi(t)\) is a monic function with integral rational functional coefficients which vanishes for \(t=\omega\), then every rational divisor of \(\Phi(t)\) has integral rational functional coefficients.  In particular the irreducible equation of \(\omega\) over the rational functional field has integral rational functional coefficients.  This proves that a functional which is fractional in \(\Omega^*\) cannot become integral merely by considering it in a higher field.

Third, every functional \(\omega\) can be made integral by multiplication by a rational integer.  If \(\omega\) satisfies a monic equation with rational functional coefficients, multiply by a common denominator of the absolute values of the coefficients.  Then a suitable multiple \(a\omega\) satisfies a monic integral equation.  Hence every functional is a quotient of two integral functionals, and the denominator may be chosen to be a natural number.

A useful criterion is the following.  Suppose there are \(m\) quantities \(\omega_1,
\ldots,
\omega_m\), not all zero, such that all products \(\omega\omega_i\) can be expressed as
\begin{equation}
 \omega\omega_i=A_{1i}\omega_1+A_{2i}\omega_2+\cdots+A_{mi}\omega_m, \tag{1}
\end{equation}
where the \(A_{ji}\) are integral rational functionals.  Then \(\omega\) is integral.  For the determinant of this linear system gives a monic equation for \(\omega\) with integral rational functional coefficients.

Consequently any integral function in integral functionals is integral.  In particular, by addition, subtraction, and multiplication one obtains only integral functionals from integral functionals.  Conjugates of an integral functional are integral.  The absolute norm of an integral functional is a natural number.

Finally, if a functional satisfies a monic equation whose coefficients are integral functionals of \(\Omega^*\), then the functional itself is integral.  To prove this, form the norm of the polynomial in a new variable \(t\).  The norm has integral rational functional coefficients and vanishes at the same functional; the preceding criterion applies.

If an integral functional is in fact a number, then it is an algebraic integer in the earlier sense.  Thus the algebraic integers are special cases of integral functionals.

\webersection{138}{Divisibility, Associated Functionals, Units}

The integral functionals obey laws of divisibility analogous to those of rational integers.  We begin with the definition.

An integral functional \(\alpha\) is called divisible by another non-zero integral functional \(\beta\) if the quotient
\[
 \gamma=\frac{\alpha}{\beta}
\]
is again an integral functional.  We then write \(\alpha=\beta\gamma\).

Two integral functionals \(\alpha,\beta\) are called associated if each is divisible by the other.  Equivalently,
\[
 \alpha=\beta\varepsilon,
\]
where \(\varepsilon\) is an integral functional whose reciprocal is also integral.  Such a functional \(\varepsilon\) is called a unit.

The absolute norm of a unit is equal to one.  Conversely, if an integral functional has absolute norm one, then it is a unit.  Indeed, its reciprocal is, up to a rational unit, the product of its conjugates divided by its norm and is integral.

Every rational unit is, in the present sense, a unit.  An integral rational functional is associated with its absolute value.  Thus the sign of an ordinary integer is here generalized to the rational unit factor of a rational functional.

For any integral functional \(\alpha\), the norm \(N(\alpha)=\alpha\alpha'\cdots\) shows that the natural integer \(N_a(\alpha)\) is divisible by \(\alpha\).  Hence there are infinitely many natural integers divisible by \(\alpha\).  Among them is a least one, say \(a\).  Every rational integer divisible by \(\alpha\) is then divisible by \(a\); for the remainder upon division by \(a\) would again be a rational integer divisible by \(\alpha\), smaller than \(a\), and so must be zero.

\webersection{139}{Greatest Common Divisor}

Let \(\alpha,\beta,\ldots\) be a finite number of non-zero integral functionals in \(\Omega^*\), and let \(x,y,\ldots\) be variables not occurring in them.  Form
\begin{equation}
 \delta=\alpha x+
\beta y+
\cdots . \tag{1}
\end{equation}
This is an integral functional.  Its norm is a homogeneous integral rational function of the variables.  If \(D\) denotes the absolute norm of \(\delta\), one may write
\begin{equation}
 N(\delta)=D E(x,y,\ldots), \tag{2}
\end{equation}
where \(E\) is a rational unit.

Now substitute arbitrary rational integers \(x_0,y_0,
\ldots\) for the variables and put
\[
 \delta_0=\alpha x_0+
\beta y_0+
\cdots .
\]
The quotient \(\delta_0/D\) is an integral functional.  This is seen by considering
\[
 \delta t-
\delta_0=\alpha(xt-x_0)+
\beta(yt-y_0)+
\cdots
\]
and its norm, ordered according to powers of \(t\).  The resulting monic equation has integral rational functional coefficients and is satisfied by \(\delta_0/D\).  Thus \(\delta_0\) is divisible by \(D\).

Since the integers \(x_0,y_0,
\ldots\) are arbitrary, it follows in particular that \(\alpha,
\beta,
\ldots\) are all divisible by \(D\).  Conversely, every common divisor of \(\alpha,
\beta,
\ldots\) divides every expression of the form (1), and hence divides \(D\).  Therefore \(D\) has the characteristic property of a greatest common divisor.

Thus every finite system of integral functionals has a greatest common divisor, determined up to multiplication by a unit.  It may be represented by a linear expression in the functionals with integral functional coefficients, in the same way as for rational integers.

\webersection{140}{Prime Functionals in the Field \(\Omega\)}

An integral functional \(\pi\), different from zero and not a unit, is called a prime functional if from a divisibility
\[
 \pi\mid \alpha\beta
\]
it follows that \(\pi\mid\alpha\) or \(\pi\mid\beta\).  This is the direct analogue of a prime number.

If \(\pi\) is a prime functional, there are natural integers divisible by \(\pi\).  Let \(p\) be the least among them.  It cannot be composite.  For if \(p=p_1p_2\) with both factors greater than one, then, since \(\pi\) is prime, one of the two factors would be divisible by \(\pi\), contradicting the minimality of \(p\).  Hence \(p\) is an ordinary rational prime.

Every natural integer divisible by \(\pi\) is divisible by this same prime \(p\).  Since
\begin{equation}
 p=\pi\alpha \tag{1}
\end{equation}
with integral \(\alpha\), taking absolute norms gives
\begin{equation}
 p^n=N_a(\pi)N_a(\alpha), \tag{2}
\end{equation}
where \(n\) is the degree of the number field.  Therefore \(N_a(\pi)\) is a power of \(p\).  Write
\begin{equation}
 N_a(\pi)=p^f. \tag{3}
\end{equation}
The integer \(f\), with \(1\leq f\leq n\), is called the degree of the prime functional \(\pi\).

\webersection{141}{Decomposition of Integral and Fractional Functionals into Prime Factors}

Every non-zero integral functional which is not a unit is divisible by a prime functional.  If the functional itself is prime, there is nothing to prove.  Otherwise it has a proper integral divisor \(\omega_1\) which is neither a unit nor associated with the original functional:
\[
 \omega=
\omega_1\alpha.
\]
Since the absolute norm of \(\alpha\) is greater than one, the absolute norm of \(\omega_1\) is strictly smaller than that of \(\omega\).  If \(\omega_1\) is still not prime, the same argument may be applied again.  The absolute norms form a strictly decreasing sequence of natural numbers, and therefore the process stops at a prime functional.

It follows that every integral functional different from zero and from the units can be decomposed into finitely many prime factors:
\begin{equation}
 \omega=\pi_1\pi_2\cdots\pi_r. \tag{1}
\end{equation}

This decomposition is unique, up to associated prime factors and the order of the factors.  The proof is the usual one.  If
\[
 \pi_1\pi_2\cdots\pi_r
\quad \text{and}\quad
 \rho_1\rho_2\cdots\rho_s
\]
are two decompositions into prime functionals, then \(\pi_1\) divides the product on the right and hence divides one of the \(\rho_j\).  Since both are prime, they are associated.  Cancelling associated factors and repeating the argument gives equality of the decompositions.

The same result applies to fractional functionals.  Since every functional is a quotient of two integral functionals, and both numerator and denominator can be decomposed into prime factors, the functional is represented uniquely as a product of prime powers with integral exponents, again up to units.

This restores, in the enlarged domain of functionals, the law of unique factorization that may fail for the algebraic integers themselves.

\webersection{142}{Integral Functions in the Field \(\Omega^*\)}

The proof of unique factorization also rests on the following analogue of Gauss's lemma.  Let
\[
 \varphi=a_0t^h+a_1t^{h-1}+\cdots+a_h,
\qquad
 \psi=b_0t^k+b_1t^{k-1}+\cdots+b_k
\]
be integral functions of one variable \(t\), whose coefficients are integral functionals not involving \(t\).  If all coefficients of the product \(\varphi\psi\) have a common prime divisor \(\pi\), then \(\pi\) divides all coefficients of \(\varphi\) or all coefficients of \(\psi\).

From this one draws an important conclusion.  A functional
\begin{equation}
 \Phi(t)=\varphi_0t^h+\varphi_1t^{h-1}+\cdots+
\varphi_h \tag{1}
\end{equation}
with coefficients in \(\Omega^*\) is integral only if all its coefficients \(\varphi_i\) are integral functionals.

Indeed, suppose not all coefficients are integral.  Let \(\mu\) be a common denominator, so that \(\mu\varphi_i=\alpha_i\) are integral and not all divisible by a prime factor \(\pi\) of \(\mu\).  If \(\Phi\) itself were integral, then it would satisfy a monic equation with integral rational functional coefficients.  After multiplying by a sufficiently high power of \(\mu\), one obtains an identity between integral functions in which the right-hand side has all coefficients divisible by \(\pi\), whereas the left-hand side does not, contradicting the lemma.

Conversely, a polynomial whose coefficients are integral functionals is integral.  Hence integral functions in \(\Omega^*\) are precisely those whose coefficients are integral functionals.

Every functional can be multiplied by a primitive rational function so as to become an integral function of the variables with coefficients in \(\Omega\).  If the functional is already integral, it is associated with such an integral function whose coefficients are integral numbers of \(\Omega\).  Thus integral functionals may be represented by ordinary integral functions, up to multiplication by rational units.

\webersection{143}{The Prime Factors of the Numbers of the Field \(\Omega\)}

Since the numbers of \(\Omega\) are among the functionals of \(\Omega^*\), the preceding theory shows that integral numbers of \(\Omega\) can be decomposed into prime factors.  These prime factors need not themselves be numbers; in general they are functionals.  In special cases prime functionals may be numbers, and then they are called prime numbers of the field.

All equations between functionals are ultimately identities between integral rational functions.  Hence they remain valid when the variables are replaced by other symbols.  It follows that an integral functional, a unit, and a prime functional remain respectively integral, a unit, and prime after such a replacement.

Consequently the prime factors of an integral functional may be represented so that they do not contain the variables on which the functional itself depends.  For every integral functional divides a number, in particular its absolute norm.  Its prime factors are therefore among the prime factors of such a number, and may be chosen with variables different from those occurring in the original functional.

Let
\begin{equation}
 E\omega=
\varphi(x,y,
\ldots) \tag{1}
\end{equation}
where \(E\) is a unit and \(\varphi\) an integral function in \(\Omega\).  If \(\omega\) is decomposed into prime factors chosen free of the variables \(x,y,
\ldots\), their product gives a number \(\omega_1\) associated with \(\omega\).  The quotient \(\varphi/
\omega_1\) is then an integral function.  Hence all coefficients of \(\varphi\) are divisible by \(\omega_1\), and therefore by \(\omega\).  Conversely, if all coefficients of \(\varphi\) are divisible by a factor, then so is \(\varphi\).

Thus the divisibility of an integral function by a number or by a prime functional is equivalent to the divisibility of all its coefficients.  This observation is the bridge between factorization of functionals and arithmetic of the numbers of \(\Omega\).

\booktitle{Seventeenth Section. Theory of Algebraic Number Fields}

\webersection{144}{Basis of an Algebraic Number Field; Discriminants}

Let \(\Omega=R(\theta)\) be an algebraic number field of degree \(n\), where \(\theta\) satisfies the irreducible equation
\begin{equation}
 f(t)=t^n+a_1t^{n-1}+\cdots+a_n=0. \tag{1}
\end{equation}
Every number of \(\Omega\) is uniquely expressible in the form
\begin{equation}
 \omega=c_1\omega_1+c_2\omega_2+
\cdots+c_n\omega_n, \tag{2}
\end{equation}
with rational coefficients, after a suitable system \(\omega_1,
\ldots,
\omega_n\) has been chosen.  Such a system is called a basis of the field.  The powers
\begin{equation}
 1,\theta,\theta^2,
\ldots,\theta^{n-1} \tag{3}
\end{equation}
form one basis.

Let \(\omega_{r,1},
\omega_{r,2},
\ldots,
\omega_{r,n}\) denote the conjugates of \(\omega_r\).  The square of the determinant
\begin{equation}
 \Delta(\omega_1,
\ldots,
\omega_n)=
\left(\det(\omega_{r,s})\right)^2 \tag{4}
\end{equation}
is a symmetric function of the conjugates and hence a rational number.  It is called the discriminant of the system \(\omega_1,
\ldots,
\omega_n\).

For the power basis one obtains the discriminant of the equation (1), equivalently
\begin{equation}
 \Delta(1,	heta,
\ldots,
\theta^{n-1})=(-1)^{n(n-1)/2}N(f'(\theta)). \tag{5}
\end{equation}
This number is non-zero.  If another system \(\alpha_1,
\ldots,
\alpha_n\) is obtained from a basis by a rational linear substitution with determinant \(H\), then
\begin{equation}
 \Delta(\alpha_1,
\ldots,
\alpha_n)=H^2\Delta(\omega_1,
\ldots,
\omega_n). \tag{6}
\end{equation}
Thus the discriminants of different bases differ by a rational square and have the same sign.  Conversely, any system of \(n\) numbers whose determinant is non-zero is a basis.

\webersection{145}{The Minimal Basis and the Field Discriminant}

A basis remains a basis if each of its elements is multiplied by a non-zero rational number.  Since every algebraic number can be made integral by multiplication by a rational integer, there exist bases all of whose elements are algebraic integers.  The discriminant of such an integral basis is a non-zero rational integer, and all such discriminants have the same sign.

Among their absolute values choose the least.  A basis of algebraic integers whose discriminant has this least absolute value is called a minimal basis.  The corresponding discriminant is called the discriminant of the field or the ground number of the field.

Let
\[
 \omega_1,
\ldots,
\omega_n
\]
be a minimal basis.  Any algebraic integer \(\alpha\in\Omega\) has a rational representation
\[
 \alpha=u_1\omega_1+
\cdots+u_n\omega_n.
\]
If some coefficient had a denominator, multiplying by a suitable rational integer and comparing discriminants would yield an integral basis with smaller absolute discriminant, which is impossible.  Hence all \(u_i\) are rational integers.

Thus the minimal basis has the fundamental property that every algebraic integer of \(\Omega\) is represented uniquely as
\begin{equation}
 \alpha=x_1\omega_1+x_2\omega_2+
\cdots+x_n\omega_n,
\qquad x_i\in\mathbb Z. \tag{1}
\end{equation}
Conversely, every expression of this form is an algebraic integer.  The minimal basis therefore gives the integral lattice of the field.

The same reasoning applies to integral functionals.  If \(\omega_1,
\ldots,
\omega_n\) is a minimal basis, then a functional
\[
 \alpha=u_1\omega_1+
\cdots+u_n\omega_n
\]
is integral precisely when the coefficients \(u_i\) are integral rational functionals.  This will be used in the construction of bases of functionals.

\webersection{146}{The Bases of Functionals}

Let \(\mu\) be an integral functional of \(\Omega^*\).  A basis of \(\mu\) is a system of \(n\) algebraic integers
\begin{equation}
 \alpha_1,
\alpha_2,
\ldots,
\alpha_n \tag{1}
\end{equation}
which is a basis of the field \(\Omega\), and which has the following property: all algebraic integers of \(\Omega\) divisible by \(\mu\), and no others, are obtained in the form
\begin{equation}
 \alpha=x_1\alpha_1+x_2\alpha_2+
\cdots+x_n\alpha_n,
\qquad x_i\in\mathbb Z. \tag{2}
\end{equation}

Every integral functional has such a basis.  To construct one, begin with a minimal basis
\[
 \omega_1,
\omega_2,
\ldots,
\omega_n.
\]
There are positive rational integers \(a\) such that \(a\omega_1\) is divisible by \(\mu\).  Let \(a_{11}\) be the least positive such integer and put
\[
 \alpha_1=a_{11}\omega_1.
\]
Next choose \(a_{22}\) least positive such that integers \(a_{12}\) exist with
\[
 \alpha_2=a_{12}\omega_1+a_{22}\omega_2
\]
divisible by \(\mu\).  Continuing in this way one obtains a triangular system
\begin{align}
 \alpha_1&=a_{11}\omega_1,\notag\\
 \alpha_2&=a_{12}\omega_1+a_{22}\omega_2,\notag\\
 &\ \,\vdots \tag{3}\\
 \alpha_n&=a_{1n}\omega_1+a_{2n}\omega_2+
\cdots+a_{nn}\omega_n,
\notag
\end{align}
where each diagonal coefficient \(a_{rr}\) is chosen minimal positive subject to divisibility by \(\mu\).

The system (3) is a basis of the functional \(\mu\).  Every integral combination of the \(\alpha_r\) is divisible by \(\mu\) by construction.  Conversely, if an algebraic integer divisible by \(\mu\) is expressed in the minimal basis, the triangular equations determine successively residues modulo \(a_{nn},a_{n-1,n-1},
\ldots,a_{11}\).  The minimality conditions force the required representation in the \(\alpha_r\).

A basis form of \(\mu\) is the linear form
\begin{equation}
 \alpha=x_1\alpha_1+
\cdots+x_n\alpha_n. \tag{4}
\end{equation}
Substituting rational integers for the variables gives precisely all algebraic integers divisible by \(\mu\).

If \(\mu=1\), a basis of \(\mu\) is a minimal basis of the field.  For arbitrary \(\mu\), all other bases of \(\mu\) are obtained from one by an integral linear substitution of determinant \(\pm1\).

\webersection{147}{The Absolute Norms of Functionals}

Let
\[
 \alpha_1,
\ldots,
\alpha_n
\]
be a basis of the functional \(\mu\).  These algebraic integers may be expressed integrally in terms of a minimal basis \(\omega_1,
\ldots,
\omega_n\):
\begin{equation}
 (\alpha_1,
\alpha_2,
\ldots,
\alpha_n)=A(\omega_1,
\omega_2,
\ldots,
\omega_n), \tag{1}
\end{equation}
where \(A\) is an integral matrix.  In the triangular basis constructed above, the determinant of \(A\) is
\[
 a_{11}a_{22}\cdots a_{nn}.
\]

Let
\begin{equation}
 \lambda=t_1\alpha_1+
\cdots+t_n\alpha_n \tag{2}
\end{equation}
be the basis form of \(\mu\).  Since the products \(\omega_r\alpha_s\) are all divisible by \(\mu\), they can be expressed in the basis \(\alpha_1,
\ldots,
\alpha_n\) with rational integral coefficients:
\begin{equation}
 \omega_r\alpha_s=\sum_u g_{rsu}\alpha_u. \tag{3}
\end{equation}
Multiplying the basis form and eliminating the \(\alpha_u\), one obtains the equation for \(\lambda\).  The product of its roots, that is, \(N(\lambda)\), factors as
\begin{equation}
 N(\lambda)=A\,T(t_1,
\ldots,t_n), \tag{4}
\end{equation}
where \(A=\det(a_{rs})\) and \(T\) is an integral rational function of the variables \(t_i\).

It remains to show that \(T\) is primitive.  If a rational prime \(p\) divided all its coefficients, then by elementary determinant theory one could choose integral forms in the variables, not all divisible by \(p\), making all corresponding linear combinations divisible by \(p\).  This would produce an algebraic integer divisible by \(\mu\) whose coordinates in the basis are not all divisible by \(p\), contradicting the defining property of the basis.

Therefore \(T\) is primitive.  Consequently the absolute value of the determinant \(A\) is the absolute norm of \(\mu\).  In particular, for the triangular basis (3) of \S 146,
\begin{equation}
 N_a(\mu)=a_{11}a_{22}
\cdots a_{nn}. \tag{5}
\end{equation}

If one applies to the variables of a basis form an integral linear substitution of determinant \(\pm1\), one obtains a new basis form of the same functional.  Equivalently, the transposed substitution acts on the basis elements.  The discriminant of a basis of \(\mu\) is therefore
\begin{equation}
 \Delta(\alpha_1,
\ldots,
\alpha_n)=N_a(\mu)^2\,d_\Omega, \tag{6}
\end{equation}
where \(d_\Omega\) is the discriminant of the field.  Conversely, if a system of integral numbers divisible by \(\mu\) has discriminant equal to this number, then it is a basis of \(\mu\).

\webersection{148}{Complete Residue System Modulo a Functional}

Two algebraic integers \(\xi,\eta\) are called congruent modulo a non-zero integral functional \(\mu\) if the difference \(\xi-
\eta\) is divisible by \(\mu\).  We write, in Gauss's notation,
\begin{equation}
 \xi\equiv\eta\pmod{\mu}. \tag{1}
\end{equation}
The modulus may be a functional, not merely a number.  Multiplying the modulus by a unit does not change the congruence.

If
\[
 \xi\equiv\xi'\pmod{\mu},
\qquad
 \eta\equiv\eta'\pmod{\mu},
\]
then every expression formed from the numbers by addition, subtraction, and multiplication remains congruent to the expression obtained by replacing each number by its congruent counterpart.  The usual rules of congruence therefore hold.

Let
\begin{equation}
 \alpha_1,
\ldots,
\alpha_n \tag{2}
\end{equation}
be a basis of the modulus \(\mu\), written in triangular form relative to a minimal basis \(\omega_1,
\ldots,
\omega_n\):
\begin{align}
 \alpha_1&=a_{11}\omega_1,\notag\\
 \alpha_2&=a_{12}\omega_1+a_{22}\omega_2,\notag\\
 &\ \vdots \tag{3}\\
 \alpha_n&=a_{1n}\omega_1+
\cdots+a_{nn}\omega_n.
\notag
\end{align}
Every algebraic integer \(\eta\) of the field can be written
\begin{equation}
 \eta=y_1\omega_1+y_2\omega_2+
\cdots+y_n\omega_n,
\qquad y_i\in\mathbb Z. \tag{4}
\end{equation}
Now form the system
\begin{equation}
 \xi=x_1\omega_1+x_2\omega_2+
\cdots+x_n\omega_n, \tag{5}
\end{equation}
where \(x_1\) runs through a complete residue system modulo \(a_{11}\), \(x_2\) modulo \(a_{22}\), and so on up to \(x_n\) modulo \(a_{nn}\).  The number of such \(\xi\) is
\[
 a_{11}a_{22}\cdots a_{nn}=N_a(\mu).
\]

We prove that every integer \(\eta\) is congruent modulo \(\mu\) to exactly one of these \(\xi\).  The condition that \(\eta-
\xi\) be divisible by \(\mu\) means that it can be written as an integral combination of the basis elements \(\alpha_i\).  Comparing coefficients in the minimal basis gives the triangular system
\begin{align*}
 y_1&=x_1+h_1a_{11}+h_2a_{12}+
\cdots+h_na_{1n},\\
 y_2&=x_2+h_2a_{22}+
\cdots+h_na_{2n},\\
 &\ \vdots\\
 y_n&=x_n+h_na_{nn}.
\end{align*}
The last equation determines \(x_n\) and \(h_n\) uniquely; then the preceding equation determines \(x_{n-1}\) and \(h_{n-1}\), and so on.  Thus all \(x_i\) are determined uniquely.

The system (5), or any system congruent to it, is therefore called a complete residue system modulo \(\mu\).  The number of individuals in a complete residue system modulo \(\mu\) is the absolute norm \(N_a(\mu)\).

\webersection{149}{Congruences}

From the existence, proved in the preceding section, of a complete residue system modulo a functional \(\mu\), there follows a series of propositions exactly analogous to the elementary propositions on congruences in rational number theory.

Let \(\alpha\) be an arbitrary integer of the field \(\Omega\), and let \(\mu\) be an integral functional serving as modulus.  If \(\alpha\) is relatively prime to \(\mu\), then two products \(\alpha\xi\) and \(\alpha\xi'\) are congruent modulo \(\mu\) only when \(\xi\) and \(\xi'\) are themselves congruent modulo \(\mu\).  Hence, when \(\xi\) runs through a complete residue system modulo \(\mu\), the products \(\alpha\xi\) also run through such a complete residue system.  Thus:

\begin{quote}
If \(\alpha\) is an integer of \(\Omega\) relatively prime to the modulus \(\mu\), and if \(\gamma\) is any integer of \(\Omega\), the congruence
\begin{equation}
 \alpha \xi \equiv \gamma \pmod{\mu} \tag{1}
\end{equation}
is soluble in an integer \(\xi\), and has exactly one solution when \(\xi\) is required to lie in a prescribed complete residue system modulo \(\mu\).
\end{quote}

If the modulus itself is represented by a number \(\beta\), this gives the following form.  If \(\alpha,\beta,\gamma\) are integers in \(\Omega\) and if \(\alpha,\beta\) are relatively prime, then two integers \(\xi,\eta\) can be chosen in \(\Omega\) so that
\begin{equation}
 \alpha\xi+\beta\eta=\gamma. \tag{2}
\end{equation}
In particular, for any two relatively prime integers \(\alpha,\beta\), the equation
\begin{equation}
 \alpha\xi+\beta\eta=1 \tag{3}
\end{equation}
can be satisfied by integers \(\xi,\eta\) of the field.

The same proposition extends to any finite system of integers.  If
\[
 \alpha,\beta,\gamma,\\ldots
\]
have no common divisor, then integers \(\xi,\eta,\zeta,
\ldots\) may be found such that
\begin{equation}
 \alpha\xi+\beta\eta+\gamma\zeta+\cdots=1. \tag{4}
\end{equation}
Indeed one first chooses coefficients in such a way that a certain linear combination of \(\beta,\gamma,\ldots\) is relatively prime to \(\alpha\), and then applies the preceding two-term result.  Conversely, if such an equation exists, no prime divisor can divide all of \(\alpha,\beta,\gamma,
\ldots\).

It follows more generally that if a system of equations of the form
\begin{equation}
 \varepsilon_i = \alpha\xi_i+\beta\eta_i+
\gamma\zeta_i+
\cdots \tag{5}
\end{equation}
is obtained by expanding a unit \(\varepsilon\) in a basis, then the coefficients \(\varepsilon_i\) have no common divisor.  Multiplying by suitable integers and adding gives (4).  This is the form of Bezout's theorem which will be needed below.

\webersection{150}{Fermat's Theorem}

The theory of congruences now gives consequences exactly corresponding to the theorems derived in rational arithmetic from Fermat's theorem.  Let \(\pi\) be a prime functional of degree \(f\), and let \(p\) be the rational prime divisible by \(\pi\).  The norm of \(\pi\) is then
\[
 N(\pi)=p^f.
\]
For any integer \(\alpha\) not divisible by \(\pi\), multiplication by \(\alpha\) permutes the non-zero residue classes modulo \(\pi\).  Multiplying all non-zero residues and cancelling their product, which is not divisible by \(\pi\), gives
\begin{equation}
 \alpha^{p^f-1}\equiv 1\pmod{\pi}. \tag{1}
\end{equation}
Multiplication by \(\alpha\) yields the equivalent form
\begin{equation}
 \alpha^{p^f}\equiv \alpha\pmod{\pi}, \tag{2}
\end{equation}
which remains valid even when \(\alpha\) is divisible by \(\pi\).  These formulae are the analogue, for algebraic number fields, of Fermat's theorem.

A polynomial congruence modulo a prime functional has no more roots than its degree.  If
\begin{equation}
 f(t)\equiv 0\pmod{\pi} \tag{3}
\end{equation}
where \(f(t)\) is an integral polynomial of degree \(m\), then there are at most \(m\) incongruent integers \(t\) modulo \(\pi\) satisfying it.  For if \(a\) is one root, then
\[
 f(t)=(t-a)f_1(t)+f(a),
\]
and all other roots either are congruent to \(a\) or satisfy a congruence of degree \(m-1\).  Induction proves the assertion.

Every integer not divisible by \(\pi\) satisfies
\begin{equation}
 \omega^{p^f-1}\equiv 1\pmod{\pi}. \tag{4}
\end{equation}
If \(a\) is the least positive exponent for which
\[
 \omega^a\equiv1\pmod{\pi},
\]
then every exponent \(l\) with \(\omega^l\equiv1\pmod{\pi}\) is a multiple of \(a\); hence \(a\mid(p^f-1)\).  We say that \(\omega\) belongs to the exponent \(a\).  If \(\omega\) belongs to \(a\), then
\[
 1,\omega,\omega^2,
\ldots,\omega^{a-1}
\]
are incongruent and give all roots of \(t^a\equiv1\pmod{\pi}\).  Among them, \(\omega^l\) belongs to the exponent \(a\) precisely when \(l\) is relatively prime to \(a\).  Therefore the number of elements belonging to exponent \(a\) is \(\varphi(a)\), whenever such elements exist.

As in rational arithmetic one proves that for every divisor \(a\) of \(p^f-1\) there are elements belonging to exponent \(a\).  The elements belonging to the greatest exponent \(p^f-1\) are called primitive roots of \(\pi\).  If \(\gamma\) is such a primitive root, then
\[
 1,\gamma,
\gamma^2,
\ldots,
\gamma^{p^f-2}
\]
form a complete system of the non-zero residues modulo \(\pi\).

Finally, if an integer \(\alpha\) is congruent modulo \(\pi\) to a rational integer, then by the ordinary Fermat theorem it satisfies
\begin{equation}
 \alpha^p\equiv \alpha\pmod{\pi}. \tag{5}
\end{equation}
Conversely, the congruence (5) characterizes precisely those residue classes which are represented by rational integers.  For a polynomial \(\psi(x,y,
\ldots)\) with integral algebraic coefficients, the congruence
\begin{equation}
 \psi(x,y,
\ldots)^p\equiv \psi(x^p,y^p,
\ldots)\pmod{\pi} \tag{6}
\end{equation}
is therefore the necessary and sufficient condition that all coefficients of \(\psi\) be congruent modulo \(\pi\) to rational integers.

\webersection{151}{Dedekind Ideals}

Dedekind based the theory of algebraic numbers on the concept of the ideal.  We now show that the theory of ideals is, in essence, the same as the theory of functionals, and describe the passage from one language to the other.

Let \(\mathfrak o\) denote, as before, the system of all integers of the algebraic number field \(\Omega\).  A system \(\mathfrak a\) of numbers contained in \(\mathfrak o\) is called an ideal if it satisfies the following two conditions:
\begin{enumerate}[label=\Roman*.]
\item The sum and difference of any two numbers of \(\mathfrak a\) again belong to \(\mathfrak a\).
\item The product of any number of \(\mathfrak a\) by any number of \(\mathfrak o\) again belongs to \(\mathfrak a\).
\end{enumerate}
The system consisting only of zero satisfies these conditions, but for simplicity is not counted as an ideal here.  The system \(\mathfrak o\) itself is an ideal.  Also, the set of all multiples in \(\mathfrak o\) of a fixed integer \(\alpha\) is an ideal; such an ideal is called a principal ideal.

The product \(\mathfrak a\mathfrak b\) of two ideals is the set of all finite sums of products \(\sum \alpha\beta\), with \(\alpha\in\mathfrak a\), \(\beta\in\mathfrak b\).  This is again an ideal.  With respect to this multiplication, \(\mathfrak o\) is the identity.

Ideals and functionals correspond as follows.  To each integral functional \(\varphi\) there belongs the ideal consisting of all integers divisible by \(\varphi\).  Associated functionals give the same ideal.  Conversely, to each ideal there correspond infinitely many integral functionals, but they are all associated.  Thus the ideal records the divisor common to all its numbers.

If the functional \(\varphi\) corresponds to \(\mathfrak a\) and \(\psi\) to \(\mathfrak b\), then the product \(\varphi\psi\) corresponds to the ideal product \(\mathfrak a\mathfrak b\).  One inclusion is immediate, since every product of an element divisible by \(\varphi\) and an element divisible by \(\psi\) is divisible by \(\varphi\psi\).  For the converse, represent \(\varphi\) and \(\psi\) as greatest common divisors of pairs of integers; expanding the products shows that every integer divisible by \(\varphi\psi\) belongs to \(\mathfrak a\mathfrak b\).

If a number or functional is divisible by the functionals belonging to an ideal, we say that it is divisible by the ideal.  Thus a factorization of a functional can also be written in the form
\begin{equation}
 \varphi=\mathfrak a\mathfrak b\cdots, \tag{1}
\end{equation}
where unit factors are suppressed.

A basis of the functional is also a basis of the corresponding ideal.  The absolute norm of the representing functional is the number which Dedekind calls the norm of the ideal.  We therefore write, if \(\varphi\) belongs to \(\mathfrak a\),
\begin{equation}
 N_a(\varphi)=N(\mathfrak a). \tag{2}
\end{equation}
The norm of an ideal is always a natural number.  If \(\mathfrak p\) is a prime ideal and \(p\) the rational prime divisible by it, then
\begin{equation}
 N(\mathfrak p)=p^f, \tag{3}
\end{equation}
and \(f\) is called the degree of the prime ideal.  In congruences the corresponding ideals may be used as moduli in place of the functionals.

Kummer's ideal numbers may be understood through the same correspondence.  If \(\alpha,\beta\) have greatest common divisor \(\chi\), and if \(\beta=\chi\varphi\), the congruence \(\alpha x\equiv0\pmod{\beta}\) is governed precisely by the functional \(\varphi\).  What Kummer called an ideal number is thereby represented by the functional or, equivalently, by the ideal attached to it.  This removes the ambiguity from the manipulation of ideal factors.

\webersection{152}{Equivalence}

We now call two functionals associated, even when fractional, if they differ only by a unit factor.  We make the following definition.

\begin{quote}
Two integral or fractional functionals \(\varphi,\psi\) in the field \(\Omega\) are called equivalent if their quotient \(\varphi:
\psi\) is associated with a number.
\end{quote}
Equivalently, there are a unit \(\varepsilon\) and a number \(\alpha\) of \(\Omega\) such that
\begin{equation}
 \varphi=\alpha\varepsilon\psi. \tag{1}
\end{equation}
If two functionals are each equivalent to a third, they are equivalent to each other.  The functionals of \(\Omega\) therefore split into classes: two functionals lie in the same class exactly when they are equivalent.  Each class is determined by any one of its representatives.

Associated functionals are equivalent, hence occur in the same class.  A class contains not merely single functionals, but all ideals determined by those functionals.  For this reason the classes are called functional classes, or more commonly ideal classes.

The collection of all integral and fractional numbers of the field, together with all units and all products of numbers by units, forms a single class.  This is called the principal class.  As an ideal class it contains the ideal \(\mathfrak o\); in what follows it is denoted by \(0\).

Every ideal class contains integral functionals.  For if \(\varphi\) is any functional and \(\alpha\) is a number chosen so that \(\alpha\varphi\) is integral, then \(\alpha\varphi\) is equivalent to \(\varphi\).

More strongly, from every ideal class one can choose a representative which is integral and relatively prime to any prescribed integral functional.  Let \(\varphi\) be an integral representative of a class \(C\), choose a number \(\alpha\) divisible by \(\varphi\), and write
\begin{equation}
 \varphi\chi=\alpha. \tag{2}
\end{equation}
Choose a number \(\beta\), divisible by \(\chi\), so that \(\beta:
\chi\) is relatively prime to the prescribed functional.  Then
\begin{equation}
 \psi=\beta:
\chi \tag{3}
\end{equation}
is equivalent to \(\varphi\) and is a representative of \(C\) relatively prime to the given functional.

\webersection{153}{The Class Number of the Field \(\Omega\)}

We now prove the important theorem:

\begin{quote}
The number of ideal classes of a field \(\Omega\) is finite.
\end{quote}
It is equivalent to prove that in every class there are integral functionals whose absolute norm does not exceed a fixed finite number depending only on the field.  Indeed every integral functional divides its absolute norm, and any fixed natural number has only finitely many non-associated divisors in the field.  Hence a uniform bound for the norm gives only finitely many classes.

Let
\[
 \omega_1,
\omega_2,
\ldots,
\omega_n
\]
be a minimal basis of \(\Omega\).  Every integer has the form
\begin{equation}
 \omega=x_1\omega_1+x_2\omega_2+
\cdots+x_n\omega_n,
\qquad x_i\in\mathbb Z. \tag{1}
\end{equation}
Choose a positive integer \(k\) and restrict the coefficients \(x_i\) to lie in the interval from \(0\) to \(k\).  Then the absolute value of each conjugate of \(\omega\) is bounded by a constant times \(k\), and the absolute norm is bounded by
\begin{equation}
 |N_a(\omega)|<Bk^n, \tag{2}
\end{equation}
where \(B\) depends only on the basis and the field.

Let \(\mu\) be an integral functional and choose \(k\) so that
\begin{equation}
 k^n<N_a(\mu)<(k+1)^n. \tag{3}
\end{equation}
The \((k+1)^n\) numbers obtained from (1) by taking \(x_i=0,1,
\ldots,k\) are more numerous than the residue classes modulo \(\mu\), whose number is \(N_a(\mu)\).  Therefore two of them are congruent modulo \(\mu\).  Their difference
\begin{equation}
 \alpha=a_1\omega_1+
\cdots+a_n\omega_n \tag{4}
\end{equation}
is non-zero, divisible by \(\mu\), and has coefficients \(|a_i|\le k\).  From (2) and (3) it follows that
\begin{equation}
 |N_a(\alpha)|<B k^n < B N_a(\mu). \tag{5}
\end{equation}
Since \(\alpha\) is divisible by \(\mu\), write
\begin{equation}
 \alpha=\mu\varphi,
\qquad N_a(\alpha)=N_a(\mu)N_a(\varphi). \tag{6}
\end{equation}
Then
\begin{equation}
 |N_a(\varphi)|<B. \tag{7}
\end{equation}

If \(\psi\) is any representative of a given class \(C\), choose \(\mu\) so that \(\mu\psi\) is a number.  Applying the preceding construction and writing \(\alpha=\mu\varphi\), one obtains a functional \(\varphi\) equivalent to \(\psi\), hence still representing \(C\), and satisfying (7).  Thus every class has a representative whose absolute norm is bounded by the same constant.  The number of classes is therefore finite.  This finite number is called the class number of \(\Omega\).

In the simplest case, where the class number is one, every functional is associated with a number.  Then every integer of the field admits a factorization into prime number factors, and the arithmetic of the field essentially agrees with rational number theory.  The field of rational numbers, the Gaussian imaginary numbers, and the field generated by the third roots of unity are among such examples.

\webersection{154}{The Group of Ideal Classes}

The ideal classes may be composed.  If \(\varphi,\psi,\varphi_1,\psi_1\) are functionals, with \(\varphi\) equivalent to \(\varphi_1\) and \(\psi\) equivalent to \(\psi_1\), then \(\varphi\psi\) is equivalent to \(\varphi_1\psi_1\).  This follows immediately from the definition, since the quotient
\[
 {\varphi\psi\over \varphi_1\psi_1}
\]
is associated with a number.

Conversely, if \(\varphi\) is equivalent to \(\varphi_1\) and \(\varphi\psi\) is equivalent to \(\varphi_1\psi_1\), then \(\psi\) is equivalent to \(\psi_1\).

Now let \(A\) and \(B\) be two ideal classes, not necessarily distinct.  Choose a representative \(\varphi\) of \(A\) and a representative \(\psi\) of \(B\).  The class \(C\) containing the product \(\varphi\psi\) is independent of the choice of representatives.  It is called the composition of \(A\) and \(B\), and is written
\begin{equation}
 C=AB=BA. \tag{1}
\end{equation}
Because this composition is derived from ordinary multiplication, it obeys the commutative and associative laws.

The identity element is the principal class \(0\).  To define the inverse of a class \(A\), take a representative \(\varphi\) and choose a number \(\alpha\) divisible by \(\varphi\).  If
\[
 \alpha=\varphi\chi,
\]
then \(\chi\) represents a class denoted \(A^{-1}\), and
\begin{equation}
 AA^{-1}=0. \tag{2}
\end{equation}
Thus the ideal classes form an abelian group, the ideal class group of the field.

Since the number of ideal classes is finite, this group is finite.  If \(h\) is its order, then for every class \(A\)
\begin{equation}
 A^h=0. \tag{3}
\end{equation}
If \(h_A\) is the least positive integer for which \(A^{h_A}=0\), then \(h_A\) divides \(h\).  Hence every functional \(\varphi\) belongs to a definite exponent dividing the class number: some power \(\varphi^{h_A}\) is associated with a number of the field.

\end{document}
